\section{Introduction}
\codeword{caseSetup} was built over a couple years for the Ontario Tech Racing team. This process is a mix of experiences gained from FSAE related activities and professional activities. The first iteration of \codeword{caseSetup} was born in the fall of 2022. The drive behind \codeword{caseSetup} was to improve the consistency of runs from user to user while reducing the chances of mistakes. The main idea behind it was to use pre-defined templates which mildly restricts the user from making un-intended modifications to the setup. This allows for consistent meshing, solving, and post processing regardless of user. 

This document here is the technical user guide, while some Computational Fluid Dynamics (CFD) theories may be included, it is not meant to be an exhaustive explination or discussion of the physics. Please refer to specialized CFD books and resources if more information is required.



\section{Directory Structure}
This process requires a specific directory structure for your cases to work. It looks like the following:
\begin{itemize}
    \item \codeword{JOB-DIRECTORY}
    \begin{itemize}
        \item 01\_incoming(YYYY-MM-DD)
        \begin{itemize}
            \item 2025-09-11\_New\_Geometry
        \end{itemize}
        \item 02\_references
        \begin{itemize}
            \item ANSA
            \begin{itemize}
                \item trial009.ansa.gz
            \end{itemize}
            \item MSH
            \begin{itemize}
                \item trial009.obj.gz
                \item fw-v1-1.obj.gz
            \end{itemize}
        \end{itemize}
        \item 03\_reports
        \begin{itemize}
            \item 051\_064\_063\_report\_2025-11-16.pptx
        \end{itemize}
        \item CASES
        \begin{itemize}
            \item 001
            \item 002
            \item 003
            \item \dots
            \item 009
            \item \begin{itemize}
                \item \textbf{caseSetup}
                \item constant
                \item system
            \end{itemize}
        \end{itemize}
    \end{itemize}
\end{itemize}

The definitions for each folder is described here:
\begin{itemize}
    \item \codeword{01\_incoming(YYYY-MM-DD)} - Contains all the incoming geometry from other departments or other files which are required for the simulation. Each instance of CAD delivered to the department must be in a dated folder with a short descriptive name.
    \item \codeword{02\_references} - Contains all the geometry files that are used in the simulation. Sub-directories are ANSA and MSH
    \item \codeword{ANSA} - Contains all the ANSA files that are used in the simulation, ideally contains the CAD files within them and not just the tri-surfaces. Should be saved as ".gz" for smaller file size.
    \item \codeword{MSH} - Contains all the tri-surfaces that are used in the simulation. Ideally this is where the original copies of all the STLs or OBJ files live. Each trial folder contains only a symbolic link referencing to this directory to reduce the number and size of geometry files.
    \item \codeword{03_reports} - Contains all the reports placed by the \codeword{pptGeneration} process or placed by the user manually.
    \item \codeword{CASES} - Contains all the simulation runs, this is where each case/trial folder lives, each trial is denoted by a 3 digit number XXX.
    \item \codeword{caseSetup} - This is the file that controls the simulations and the outputs.
    \item \codeword{system} - This directory contains the files required for the simulation to run.
    \item \codeword{constant} - Like \codeword{system} this directory also contains the files required for the simulation and also in \codeword{constant/triSurfaces} is where the geometry for the simulation is stored for this case. It is ideally a symbolic link so that repeated geometries are not copied many times in many cases.
\end{itemize} 

Once all these directories are setup, we are ready to explore the rest of the features.


